\documentclass[10pt,a4paper]{amsart}
\usepackage[margin=19mm]{geometry}
\usepackage{amsmath,amssymb,mathtools}
\usepackage[scr=rsfs,cal=euler]{mathalfa}
\usepackage{xfrac}
\usepackage[unicode,hidelinks,bookmarks=false]{hyperref}
\newcommand{\ie}{i.e.\ }
\newcommand{\cf}{cf.\ }
\newcommand{\resp}{respectively\ }
\newcommand{\Subsection}{\subsection}
\providecommand{\nobreakdash}{\nobreak\hbox{-}}
\setlength{\emergencystretch}{2em}
\multlinegap0pt
\usepackage{bm}
\usepackage{bbm}
\usepackage{dsfont}
\usepackage{xspace}
\usepackage{cancel}
\usepackage{mathtools}
\usepackage{amsthm}
\makeatletter
\@ifundefined{claim}{\newtheorem*{claim}{Claim}}{}
\@ifundefined{thm}{\newtheorem{thm}{Theorem}}{}
\makeatother
\title[Numerical Fredholm determinants for matrix-valued kernels on]{Numerical Fredholm determinants for matrix-valued kernels on the real line}
\author{Erika Gallo, John Zweck, Yuri Latushkin}
\date{Source: arXiv:2507.22875v1 (CC BY 4.0)}
\begin{document}
\maketitle

\noindent\emph{Source and licence.} Numerical Fredholm determinants for matrix-valued kernels on the real line. Version arXiv:2507.22875v1 (CC BY 4.0), distributed under the Creative Commons Attribution 4.0 International licence, as recorded in the arXiv licence metadata for this exact version and shown on the abstract page https://arxiv.org/abs/2507.22875v1. Full rights notice: licenses/arxiv-2507.22875-CC-BY-4.0.txt.\par\smallskip

\noindent\textit{Editorial scope.} Counted unit: the Thm whose source label is \texttt{thmcompbound} in the article (source0.tex lines 1424--1429, source label \texttt{thmcompbound}), delivered together with the article's own complete original proof of it (source0.tex lines 1430--1524, one \texttt{proof} environment of 95 lines). No mathematics is written, repaired or re-derived: statement and proof are the article's own text line for line. Required context retained uncounted: Qab (lines 1416--1418). Every definition, display and label of those lines is retained unchanged. Omitted from this package and not redistributed: 1--1415, 1419--1423, 1525--1806. Nothing inside a retained excerpt is omitted or altered: every retained body is delivered exactly as the article's lines are written, so no omission receipt of any kind accompanies this package. Citations: the retained excerpts carry no citation command at all, so no bibliography entry of the article is transcribed and no reference is redistributed. Reference adaptation: none was needed and none was made; every reference inside the delivered unit resolves inside it, so no unresolved reference or citation is produced. Printed slips kept literal: no printed word, symbol, hypothesis or number of the counted statement or of its proof was altered. Layout and class: the article's own class is \texttt{birkjour} with packages including verbatim, amsmath, amssymb, amsthm, cancel, hyperref; the wrapper is the lane's pinned \texttt{amsart} editorial wrapper at 10pt on A4, which restates the theorem-like environments the retained text needs and adds the article's own macro definitions that the retained text needs (none), with extra wrapper packages (bm, bbm, dsfont, xspace, cancel, mathtools). The delivered numbering is the unit's own. Assets: no retained span contains a figure environment, a graphics-inclusion command, a PDF-page inclusion, a table or a TikZ picture; no image file is copied into this package, the package declares no asset dependency and no image file is redistributed with it. Licence: version arXiv:2507.22875v1 of this article is distributed under the Creative Commons Attribution 4.0 International licence, as recorded in the arXiv licence metadata for this exact version and shown on the abstract page https://arxiv.org/abs/2507.22875v1. A full rights notice is delivered at licenses/arxiv-2507.22875-CC-BY-4.0.txt. Characters: every retained excerpt is pure ASCII, so the delivered engine, XeLaTeX with its default Latin Modern text font, reports no missing character.
\begin{equation}\label{Qab}
    Q_{[a,b]}(f) = \sum_{k=1}^{2M+1} w_k f(x_k),
\end{equation}

\begin{thm}\label{thmcompbound}
Suppose that $f \in C^{r-1,1}([a,b]^n,\mathbb C)$ for some $1\leq r \leq 4.$ Let $Q_{[a,b]}^{(n)}$ be the $n$-dimensional quadrature rule induced by the composite Simpson's quadrature rule $Q_{[a,b]}$ in (\ref{Qab}). Then
\begin{equation}
 \left| Q_{[a,b]}^{(n)}(f) - \int_{[a,b]^n} f(\mathbf{x}) d\mathbf{x}  \right| \leq c_r  4^{-r}  |f|_r n(b-a)^n  (\Delta x_{\operatorname{max}})^r.
\end{equation}
\end{thm}

\begin{proof}
Since 
%\begin{equation}
   $ [a,b]^n = \bigcup_{j_1,\dots,j_n=1}^M I_{j_1} \times \dots \times I_{j_n}$,
%\end{equation}
\begin{equation}\label{Qnfab}
    Q_{[a,b]}^{(n)}(f) = \sum_{j_1, \dots, j_n=1}^M Q_{j_1 \dots j_n}(f),
\end{equation}
where
\begin{eqnarray} \nonumber 
    Q_{j_1 \dots j_n}(f) &=& \sum_{k_1, \dots, k_n = -1}^{1} w_{j_1,k_1} \dots w_{j_n,k_n} f(x_{2j_1 - k_1}, \dots, x_{2j_n - k_n})\\  
    &\cong& \int_{I_{j_1}\times \dots \times I_{j_n}} f(x_1,\dots, x_n) dx_1 \dots dx_n.
\end{eqnarray}
For each $m \leq n$ let 
\begin{equation}
    F_{j_1 \dots j_m}(y_{m+1}, \dots, y_n) := \int_{I_{j_1} \times \dots \times I_{j_m}} f(x_1, \dots, x_m, y_{m+1}, \dots, y_n) dx_1 \dots dx_m
\end{equation}
be the partial iterated integral of $f,$ which is approximated by the partial quadrature rule 
\begin{align}
    &(Q_{j_1 \dots j_m} f)(y_{m+1}, \dots, y_n)  
    \nonumber
    \\
    &= \sum_{k_1, \dots, k_m = -1}^1  w_{k_1}\dots w_{k_m} f(x_{2j_1 - k_1}, \dots, x_{2j_m - k_m}, y_{m+1}, \dots, y_n).  
\end{align}
The identity $Q_{j_1 \dots j_n}(f)  = Q_{j_n} (Q_{j_1 \dots j_{n-1}} f)$
enables us to use  induction to prove the following claim. 

\begin{claim}
\begin{equation}\label{anclaim}
    \left|  Q_{j_1 \dots j_n}(f) - \int\limits_{I_{j_1} \times \dots \times I_{j_n}}  f(\mathbf x) d\mathbf{x}   \right| 
     \,\,\leq\,\,
     c_r 4^{-r}|f|_r \,n     (\Delta x_{\operatorname{max}})^{r}  
     \prod\limits_{k=1}^n \Delta x_{j_k}
\end{equation}
\end{claim}

\begin{proof} 
We calculate that 
\begin{align} 
&\left| Q_{j_1 \dots j_n}(f) \,\,-\,\, \int_{I_{j_1 \times \dots \times I_{j_n}}}   f(\mathbf{x}) d\mathbf{x}  \right| 
 \nonumber
\\
\,\,=\,\,& \left| 
Q_{j_n}(Q_{j_1,\dots,j_{n-1}} f) \,\,-\,\,  \int_{I_{j_n}}\int_{I_{j_1} \times \dots \times I_{j_{n-1}}}  f(\mathbf{x}, x_n) d\mathbf{x} \,dx_n  \right|
 \nonumber
\\
\,\,\leq\,\, & \left| Q_{I_{j_n}} (Q_{j_1,\dots,j_{n-1}} f) \,\,-\,\, \int_{I_{j_n}} (Q_{j_1 \dots j_{n-1}} f)(x_n) dx_n  \right|
 \nonumber 
\\
&\,\,+\,\,   \int_{I_{j_n}} \left| (Q_{j_1,\dots,j_{n-1}} f)(x_n)\,\, -\,\, 
\int_{I_{j_1} \times \dots \times I_{j_{n-1}}}  f(\mathbf{x}, x_n) d\mathbf{x}  \right| dx_n
 \nonumber 
\\  
\,\,\leq\,\, & c_r 4^{-r} (\Delta x_{\operatorname{max}})^{r}  
\left( | Q_{j_1 \dots j_{n-1}} f |_r \Delta x_{j_n}  
+  |f|_r (n-1)     \prod\limits_{k=1}^{n} \Delta x_k   \right)
\nonumber 
\\
\,\,\leq\,\,&   c_r 4^{-r}|f|_r \,n     (\Delta x_{\operatorname{max}})^{r}  
     \prod\limits_{k=1}^n \Delta x_{j_k},
     \nonumber
 \end{align}
since 
\begin{align}
&\left| (Q_{j_1 \dots j_{n-1}} f)^{(r)}(x_n) \right| 
\nonumber
\\
\,\,\leq\,\,&
\sum_{k_1,\dots,k_{n-1} = -1}^1 \!\!\!\!\! \!\!\! w_{j_1,k_1} \dots w_{j_{n-1},k_{n-1}} \left| \partial_n^r f(x_{2j_1 - k_1}, \dots, x_{2j_{n-1} - k_{n-1}}, x_n) \right|
\nonumber 
\\ 
\,\,\leq\,\,&  |f|_r  \prod\limits_{k=1}^{n-1}(w_{j_k,-1} + w_{j_k,0} + w_{j_k,1})
\,\,\leq\,\,  |f|_r  \prod\limits_{k=1}^{n-1} \Delta x_{j_k}.
\nonumber
\end{align}
\end{proof}
To complete the proof of the Theorem, by (\ref{Qnfab}) and (\ref{anclaim}), we have
\begin{align}
\left | Q_{[a,b]}^{(n)}f  - \int_{[a,b]^n}  f(\mathbf{x}) d\mathbf{x} \right|
&\,\,\leq\,\, 
\sum\limits_{j_1 \dots j_n = 1}^M \left| 
Q_{j_1 \dots j_n} f  \,\,- 
\int\limits_{I_{j_1}\times \dots \times I_{j_n}} f(\mathbf{x}) d\mathbf{x} 
\right|
\nonumber
\\
&\,\,\leq\,\, 
c_r 4^{-r}|f|_r \,n     (\Delta x_{\operatorname{max}})^{r}  
   \sum\limits_{j_1 \dots j_n = 1}^M   \prod\limits_{k=1}^n \Delta x_{j_k},
\end{align}
where
$ \sum\limits_{j_1 \dots j_n = 1}^M   \prod\limits_{k=1}^n \Delta x_{j_k} =
(b-a)  \sum\limits_{j_1 \dots j_{n-1} = 1}^M   \prod\limits_{k=1}^{n-1} \Delta x_{j_k}
= (b-a)^n$.
\end{proof}

\end{document}
